% SPDX-License-Identifier: CC-BY-SA-4.0
% Author: Matthieu Perrin
% Part: <Nom de la partie>
% Section: <Nom de la section>
% Sub-section: <Nom de la sous-section>  % (facultatif, laisser vide si non utilisé)
% Frame: <Titre de la slide>

\begingroup

\begin{frame}{Problème des matrices mortelles}

  \on[text, top=-4mm]{
    \Probleme{MortalMatrix}{
      Un ensemble fini $\{M_1, ..., M_k\}$ de matrices carrées
      \begin{itemize}
      \item Les coefficients sont entiers
      \item Les matrices de même dimension $n \times n$ ($n \ge 3$)
      \end{itemize}
    }{
      Existe-t-il un produit fini qui donne la matrice nulle ?
      \begin{itemize}
      \item on peut utiliser plusieurs fois la même matrice
      \end{itemize}
    }
  }

  \onBlock[y=6mm, anchor=north]{Théorème (admis)}{
    \textsc{MortalMatrix} est indécidable.
  }

  \onExampleBlock[y=-7mm, anchor=north]{Exemples}{
    \begin{itemize}
    \item {\footnotesize$\left\{
      \begin{pmatrix}
        0 & 1 & 0 \\
        0 & 0 & 1 \\
        0 & 0 & 0
      \end{pmatrix}
      \right\}$}
      est une instance positive
    \item {\footnotesize$\left\{
      \begin{pmatrix}
        1 & 0 & 0 \\
        0 & 1 & 0 \\
        0 & 0 & 1
      \end{pmatrix}
      \right\}$}
      est une instance négative
    \end{itemize}
  }

  \footnoteref{M. Paterson. \emph{Unsolvability in 3×3 matrices}. Studies in Applied Mathematics (1970)}
  
\end{frame}

\endgroup
